\begin{solution}{89}
\begin{subsolution}
若
\begin{equation}
k A _ {0} \leq \mu m g \quad \text {或} \quad 0 <   A _ {0} \leq \frac {\mu m g}{k}
\label{eq:1.s89}
\end{equation}
小球静止于其初始位置
\begin{equation}
x = - A _ {0} <   0
\label{eq:2.s89}
\end{equation}
若
\begin{equation}
A _ {0} > \frac {\mu m g}{k} \qquad \text {或} \qquad k A _ {0} > \mu m g
\label{eq:3.s89}
\end{equation}
小球能向右运动。设小球第一次向右运动到速度为零时弹簧的伸长量为 $x$ （其符号暂未确定），根据功能原理有
\begin{equation}
\frac {1}{2} k A _ {0} ^ {2} - \frac {1}{2} k x ^ {2} = \mu m g (A _ {0} + x)
\label{eq:4.s89}
\end{equation}
由此得
\begin{equation}
x = A _ {0} - \frac {2 \mu m g}{k}
\label{eq:5.s89}
\end{equation}
当 $x \leq 0$ ，即
\begin{equation}
\frac {\mu m g}{k} <   A _ {0} \leq \frac {2 \mu m g}{k}
\label{eq:6.s89}
\end{equation}
弹簧处于压缩或自由状态，且小球所受向右的弹力
\begin{equation}
F = k \left| B _ {1} \right| = 2 \mu m g - k A _ {0} \leq \mu m g
\label{eq:7.s89}
\end{equation}
故小球最终静止于原点左边或原点
\begin{equation}
x = A _ {0} - \frac {2 \mu m g}{k} \leq 0
\label{eq:8.s89}
\end{equation}
\eqref{eq:6.s89}式中取等号时，\eqref{eq:8.s89}式也取等号。

当 $x > 0$ ，即
\begin{equation}
A _ {0} > \frac {2 \mu m g}{k}
\label{eq:9.s89}
\end{equation}
弹簧处于伸长状态，小球所受向左的弹力
\begin{equation}
k x = k A _ {0} - 2 \mu m g
\label{eq:10.s89}
\end{equation}
若 $kx \leq \mu mg$ ，即
\begin{equation}
\frac {2 \mu m g}{k} <   A _ {0} \leq \frac {3 \mu m g}{k}
\label{eq:11.s89}
\end{equation}
小球最终静止于原点右边
\begin{equation}
x = B _ {1} = A _ {0} - \frac {2 \mu m g}{k} > 0
\label{eq:12.s89}
\end{equation}
\end{subsolution}
\begin{subsolution}
解法 (一)

设在小球完成第一次向右运动后的瞬间，记弹簧的伸长为 $B_{1}$ 。按题意，此后小球至多只能向左运动。若
\begin{equation}
B _ {1} > \frac {\mu m g}{k} \qquad \text {或} \qquad A _ {0} > \frac {3 \mu m g}{k}
\label{eq:13.s89}
\end{equation}
小球能向左运动。设小球第一次向左运动到速度为零时弹簧的伸长量为 $x$ （其符号暂未确定），根据功能原理有
\begin{equation}
\frac {1}{2} k B _ {1} ^ {2} - \frac {1}{2} k x ^ {2} = \mu m g (B _ {1} - x)
\label{eq:14.s89}
\end{equation}
由此得
\begin{equation}
x = \frac {2 \mu m g}{k} - B _ {1} = \frac {4 \mu m g}{k} - A _ {0}
\label{eq:15.s89}
\end{equation}
当 $x \geq 0$ ，即
\begin{equation}
\frac {3 \mu m g}{k} <   A _ {0} \leq \frac {4 \mu m g}{k}
\label{eq:16.s89}
\end{equation}
弹簧处于伸长或自由状态，小球所受向左的弹力
\begin{equation}
k x = 4 \mu m g - k A _ {0} \leq \mu m g
\label{eq:17.s89}
\end{equation}
小球最终静止于原点右边或原点
\begin{equation}
x = \frac {4 \mu m g}{k} - A _ {0} \geq 0
\label{eq:18.s89}
\end{equation}
\eqref{eq:18.s89}式中取等号时，\eqref{eq:16.s89}式也取等号。

当 $x < 0$ ，即
\begin{equation}
A _ {0} > \frac {4 \mu m g}{k}
\label{eq:19.s89}
\end{equation}
弹簧处于压缩状态，且小球所受向右的弹力
\begin{equation}
k | x | = k A _ {0} - 4 \mu m g
\label{eq:20.s89}
\end{equation}
如果 $k|x| \leq \mu mg$ ，即当
\begin{equation}
\frac {4 \mu m g}{k} <   A _ {0} \leq \frac {5 \mu m g}{k}
\label{eq:21.s89}
\end{equation}
小球最终静止于
\begin{equation}
x \equiv - A _ {1} = \frac {4 \mu m g}{k} - A _ {0} <   0
\label{eq:22.s89}
\end{equation}

解法（二）

设在小球完成第一次向右运动后的瞬间，记弹簧的伸长为 $B_{1}$ 。按题意，此后小球至多只能向左运动。将将 $B_{1}$ 视为 $A_{0}$ ，利用（1）的结果\eqref{eq:1.s89}\eqref{eq:2.s89}\eqref{eq:6.s89}\eqref{eq:8.s89}以及 $B_{1}$ 的表达式\eqref{eq:12.s89}可知，当
\begin{equation}
\frac {3 \mu m g}{k} <   A _ {0} \leq \frac {4 \mu m g}{k}
\label{eq:23.s89}
\end{equation}
时，小球最终静止于
\begin{equation}
x = \frac {4 \mu m g}{k} - A _ {0} \geq 0
\label{eq:24.s89}
\end{equation}
\eqref{eq:23.s89}式取等号时，\eqref{eq:24.s89}式也取等号。当
\begin{equation}
\frac {4 \mu m g}{k} <   A _ {0} \leq \frac {5 \mu m g}{k}
\label{eq:25.s89}
\end{equation}
时，小球最终静止于
\begin{equation}
x \equiv - A _ {1} = \frac {4 \mu m g}{k} - A _ {0} <   0
\label{eq:26.s89}
\end{equation}
\end{subsolution}
\begin{subsolution}
设小球第1次向右运动至速度为零的位置为 $B_{1}$ （ $B_{1}>0$ ），第1次返回至速度为零的位置为 $-A_{1}$ （ $A_{1}>0$ ）；…；第 $n$ 次向右运动至速度为零的位置为 $B_{n}$ （ $B_{n}>0$ ），第 $n$ 次返回至速度为零的位置为 $-A_{n}$ （ $A_{n}>0$ ）。由\eqref{eq:12.s89}\eqref{eq:22.s89}式并类推有
\begin{equation}
A _ {0} - B _ {1} = \frac {2 \mu m g}{k},
\label{eq:27.s89}
\end{equation}
\begin{equation}
B _ {1} - A _ {1} = \frac {2 \mu m g}{k},
\label{eq:28.s89}
\end{equation}
\begin{equation}
A _ {1} - B _ {2} = \frac {2 \mu m g}{k},
\label{eq:29.s89}
\end{equation}
\begin{equation}
A _ {n - 1} - B _ {n} = \frac {2 \mu m g}{k},
\label{eq:30.s89}
\end{equation}
\begin{equation}
B _ {n} - A _ {n} = \frac {2 \mu m g}{k}.
\label{eq:31.s89}
\end{equation}
由此得
\begin{equation}
A _ {n} = A _ {0} - \frac {4 n \mu m g}{k}
\label{eq:32.s89}
\end{equation}
将 $A_{n}$ 视为初始压缩量 $A_{0}$ ，利用（1）（2）的结果\eqref{eq:1.s89}\eqref{eq:2.s89}\eqref{eq:6.s89}\eqref{eq:8.s89}\eqref{eq:11.s89}\eqref{eq:12.s89}\eqref{eq:23.s89}\eqref{eq:24.s89}\eqref{eq:25.s89}\eqref{eq:26.s89}式以及 $A_{n}$ 的表达式\eqref{eq:32.s89}可知，小球最终静止于
\begin{equation}
x = \frac {4 n \mu m g}{k} - A _ {0} <   0, \text {当} \frac {4 n \mu m g}{k} <   A _ {0} \leq \frac {(4 n + 1) \mu m g}{k},
\label{eq:33.s89}
\end{equation}
\begin{equation}
x = A _ {0} - \frac {2 (2 n + 1) \mu m g}{k} \leq 0, \text {当} \frac {(4 n + 1) \mu m g}{k} <   A _ {0} \leq \frac {2 (2 n + 1) \mu m g}{k},
\label{eq:34.s89}
\end{equation}
\begin{equation}
x = A _ {0} - \frac {2 (2 n + 1) \mu m g}{k} > 0, \text {当} \frac {2 (2 n + 1) \mu m g}{k} <   A _ {0} \leq \frac {(4 n + 3) \mu m g}{k},
\label{eq:35.s89}
\end{equation}
\begin{equation}
x = \frac {4 (n + 1) \mu m g}{k} - A _ {0} \geq 0, \text {当} \frac {(4 n + 3) \mu m g}{k} <   A _ {0} \leq \frac {4 (n + 1) \mu m g}{k}
\label{eq:36.s89}
\end{equation}
\eqref{eq:34.s89}\eqref{eq:36.s89}式中当后式取等号时前式也取等号。
\end{subsolution}
\begin{subsolution}
设小球在由开始运动直至静止整个过程中通过的总路程为 $s$ ，设由功能原理有
\begin{equation}
\frac {1}{2} k A _ {0} ^ {2} - \frac {1}{2} k x ^ {2} = \mu m g s
\label{eq:37.s89}
\end{equation}
式中，$x$ 的值如\eqref{eq:33.s89}\eqref{eq:34.s89}\eqref{eq:35.s89}\eqref{eq:36.s89}式所示。由此得
\begin{equation}
s = \frac {k}{2 \mu m g} (A _ {0} ^ {2} - x ^ {2})
\label{eq:38.s89}
\end{equation}
即
\begin{equation}
s = 4 n \left(A _ {0} - \frac {2 n \mu m g}{k}\right), \text {当} \frac {4 n \mu m g}{k} <   A _ {0} \leq \frac {(4 n + 1) \mu m g}{k},
\label{eq:39.s89}
\end{equation}
\begin{equation}
s = 2 (2 n + 1) \left(A _ {0} - \frac {(2 n + 1) \mu m g}{k}\right), \text {当} \frac {(4 n + 1) \mu m g}{k} <   A _ {0} \leq \frac {(4 n + 3) \mu m g}{k},
\label{eq:40.s89}
\end{equation}
\begin{equation}
s = 4 (n + 1) \left(A _ {0} - \frac {2 (n + 1) \mu m g}{k}\right), \text {当} \frac {(4 n + 3) \mu m g}{k} <   A _ {0} \leq \frac {4 (n + 1) \mu m g}{k}
\label{eq:41.s89}
\end{equation}
\end{subsolution}
\end{solution}